\chapter{Uvod}


Velik broj današnjih uređaja povezan je u neku mrežu, a velika većina njih pripada lokalnim mrežama - kućnima, poslovnima ili vojnima. Lokalne su mreže time 
postale sveprisutne, no njihova sigurnost nije pratila brzinu razvoja i pristupačnosti tih istih mreža. Tek nakon niza sigurnosnih incidenata počinje rasti svijest 
poslovnih i privatnih korisnika o potrebi za zaštitom.

Kako bi se incidenti, a posljedično i troškovi za poslovne korisnike, smanjili, na tržištu se pojavilo mnoštvo programskih rješenja 
\cite{splunk_docs}\cite{security_onion}, pa i zasebnih zanimanja poput analitičara u sigurnosno-operativnom centru (SOC). Takva su rješenja dostupna 
prije svega onima koji ih mogu platiti. Za neprekidan nadzor većim je organizacijama isplativo zaposliti tim analitičara raspoređenih u smjene. Postoje i 
automatizirana rješenja niže razine zaštite, no zajednički im je nedostatak cijena odnosno manjak resursa.

Cilj je ovog rada ponuditi programsko rješenje namijenjeno privatnim korisnicima i manjim tvrtkama, kojima visoki troškovi profesionalnih sustava predstavljaju 
prepreku koju veće organizacije lakše preskoče. U radu su razrađena i ispitana dva moguća rješenja navedenog problema. Kao temeljni alat za prikupljanje podataka 
korišten je Zeek, dok je analiza prikupljenih podataka provedena nad dvama podsustavima za pohranu - baza InfluxDB uz agenta Telegraf te baza Elasticsearch uz agenta Filebeat. 
U oba se slučaja prikupljeni podatci vizualiziraju u alatu Grafana.

Nakon opisa rješenja i postupka čitatelju se predstavljaju mogući smjerovi razvoja ovoga implementacijskog rješenja te daljnji koraci predviđeni za nastavak rada.